\begin{spoken-clean}[00:00:00 - 00:00:12]
Let me continue with the statement. We were trying to state this theorem that I informally described as \qt{cutting the potato}. But now we are trying to be more serious, and I will introduce the proper notations.
\end{spoken-clean}

\begin{math-stroke}[Transition from Spherical Coordinates to Slicing Ideas]
The previous discussion led to the volume integral for the unit ball:
\[
\mu_3(B_1) = \mu_3(\Psi(D)) = \int_{(0,1)\times(0,2\pi)\times(-\pi/2,\pi/2)} y_1^2 \cos y_3 \, dy = \, \textcolor{red}{?}
\]
This naturally introduces the concept of ``Slicing ideas'':
\begin{center}
%\includegraphics[width=0.4\textwidth]{path/to/cavalieri_principle_concept.png} % Placeholder for visible board image
\end{center}
\begin{explanation-of-steps}
The board visually summarizes the core idea of Cavalieri's principle: for an $n$-dimensional body, its total volume can be computed by summing the "areas" (measures) of its $(n-1)$-dimensional slices, multiplied by an infinitesimal thickness (i.e., $\text{volume}(A) \approx \sum A_x \, dx$).
\end{explanation-of-steps}
\end{math-stroke}

\begin{spoken-clean}[00:00:12 - 00:02:20]
This theorem, you will see, is a bit the opposite of the Change of Variables formula. The Change of Variables formula has a very simple statement but a very hard proof. This has a very complicated statement, but the proof is very easy. We have our measurable set. Let me do a drawing of it. We have our set $A$ that is contained in a bounding box. Here, we will put the $y$-axis, and this is the $x$-axis. This is actually $\mathbb{R}^k$, and on the $y$-axis, we have $\mathbb{R}^{n-k}$. The whole space is $\mathbb{R}^n$, the product of the two (i.e., $\mathbb{R}^n = \mathbb{R}^k \times \mathbb{R}^{n-k}$).
\end{spoken-clean}

\begin{math-stroke}[Geometric Visualization Setup]
\begin{center}
\begin{tikzpicture}[scale=1.0]
    % Axes
    \draw[->, thick, profgreen] (-0.5,0) -- (4,0) node[right] {$x \in \mathbb{R}^k$};
    \draw[->, thick, profblue] (0,-0.5) -- (0,3) node[above] {$y \in \mathbb{R}^{n-k}$};
    \node[above right] at (3.5, 2.5) {$\mathbb{R}^n$};

    % Bounding Box (dashed)
    \draw[dashed] (0.5,0.5) rectangle (3.5,2.5);

    % Set A
    \draw[thick] (1.5,1) to[out=45,in=135] (2.5,1.5) to[out=-45,in=45] (2,2) to[out=-135,in=-45] (1,1.5) to[out=135,in=-135] cycle;
    \node at (2, 1.5) {$A$};

    % Slice at x
    \draw[red, very thick] (1.75, 0.5) -- (1.75, 2.5);
    \node[below, xshift=-0.1cm, red] at (1.75, 0) {$x$};
    \node[right, red] at (1.75, 1.5) {$A_x$};
\end{tikzpicture}
\end{center}
\begin{explanation-of-steps}
The professor draws a set $A$ within a bounding box in $\mathbb{R}^n$, split into $\mathbb{R}^k$ ($x$-axis) and $\mathbb{R}^{n-k}$ ($y$-axis). A vertical slice at a specific $x$ value is indicated in red and labeled $A_x$. This $A_x$ represents the $y$-coordinates corresponding to a fixed $x$ within $A$.
\end{explanation-of-steps}
\end{math-stroke}

\begin{spoken-clean}[00:02:20 - 00:03:37]
We are saying that our set is contained in some very large bounding box. Now, for an $x$ inside of my box, I will consider the vertical slice. The relevant piece is this slice over here. This is what I am calling $A_x$. The set $A_x$ will be the projection onto this vertical axis. So, you take this slice and you project it.
\end{spoken-clean}

\begin{nice-box}[Cavalieri's Principle for Sets]
\begin{theorem}[Cavalieri's Principle for Sets]
Let $\mathbb{R}^n = \mathbb{R}^k \times \mathbb{R}^{n-k}$, with coordinates $(x,y) \in \mathbb{R}^k \times \mathbb{R}^{n-k}$.
Let $A \subset \mathbb{R}^n$ be a Jordan measurable and bounded set, such that $A \subset (-C,C)^n$ for some $C > 0$.

We define the slices of $A$ for each $x \in (-C,C)^k$ as:
\[
A_x = \{ y \in \mathbb{R}^{n-k} \mid (x,y) \in A \} \subset (-C,C)^{n-k}
\]
We then define the upper and lower marginal densities (inner and outer measures of the slices):
\[
\overline{\varphi}(x) = \mu_{n-k}^{\text{out}}(A_x) \quad \text{and} \quad \underline{\varphi}(x) = \mu_{n-k}^{\text{in}}(A_x)
\]
Then, both functions $\overline{\varphi}, \underline{\varphi}: [-C,C]^k \to \mathbb{R}$ are Riemann integrable.
And their integrals over $[-C,C]^k$ are equal to the $n$-dimensional measure of $A$:
\[
\int_{[-C,C]^k} \overline{\varphi}(x) \, dx = \int_{[-C,C]^k} \underline{\varphi}(x) \, dx = \mu_n(A)
\]
\end{theorem}
\end{nice-box}

\begin{spoken-clean}[00:03:37 - 00:04:08]
Let us prove that. Remember that whenever we want to prove something about Jordan measures, we approximate our set from the inside and the outside.
\end{spoken-clean}

\begin{proof}[Proof of the Cavalieri's Principle for Sets]
\begin{math-stroke}[Approximation by Dyadic Cubes (Upgraded Visual)]
\begin{center}
\begin{tikzpicture}[scale=1.8]
    % Faint Background Grid to emphasize the "pixels"
    \draw[step=0.25, gray!15, very thin] (0.4, 0.4) grid (3.6, 2.1);

    % Axes
    \draw[->, thick, profgreen] (-0.2,0) -- (4,0) node[right] {$x \in \mathbb{R}^k$};
    \draw[->, thick, profgreen] (0,-0.2) -- (0,2.5) node[above] {$y \in \mathbb{R}^{n-k}$};

    % 1. The True Smooth Set A (Drawn behind the cubes)
    \draw[thick, gray, dashed, fill=gray!10] (1.1,0.6) to[out=45,in=180] (1.8,1.6) to[out=0,in=135] (3.1,1.1) to[out=-45,in=0] (2.5,0.6) to[out=180,in=-135] cycle;
    \node[gray, font=\bfseries] at (2.9, 1.2) {$A$};

    % 2. Inner Approximation F (Yellow, perfectly blocky)
    \draw[thick, profyellow!80!black, fill=profyellow!30] 
        (1.25, 0.75) -- (1.25, 1.25) -- (1.5, 1.25) -- (1.5, 1.5) -- (2.25, 1.5) -- (2.25, 1.25) -- (2.5, 1.25) -- (2.5, 1.0) -- (2.75, 1.0) -- (2.75, 0.75) -- cycle;
    \node[profyellow!80!black] at (2.0, 1.1) {$F$};

    % 3. Outer Approximation G (Red Outline, perfectly blocky)
    \draw[very thick, profred] 
        (1.0, 0.5) -- (1.0, 1.5) -- (1.25, 1.5) -- (1.25, 1.75) -- (2.5, 1.75) -- (2.5, 1.5) -- (3.0, 1.5) -- (3.0, 1.25) -- (3.25, 1.25) -- (3.25, 0.5) -- cycle;
    \node[profred] at (3.1, 1.6) {$G$};

    % 4. The vertical slice at x
    \draw[thick, profblue] (1.625, 0) -- (1.625, 2.2);
    \node[below, profblue] at (1.625, 0) {$x$};
    
    % Highlighting the trapped 1D slice bounds on the slice line
    \draw[line width=2.5pt, profred] (1.625, 0.5) -- (1.625, 1.75); % Outer slice G_x
    \draw[line width=2.5pt, gray] (1.625, 0.68) -- (1.625, 1.55); % True slice A_x (approx bound)
    \draw[line width=2.5pt, profyellow] (1.625, 0.75) -- (1.625, 1.5); % Inner slice F_x
    
    % Slice Labels
    \node[profred, right] at (1.625, 1.85) {$G_x$};
    \node[profyellow!80!black, right] at (1.625, 0.85) {$F_x$};
\end{tikzpicture}
\end{center}
\begin{explanation-of-steps}
To rigorously prove Cavalieri's Principle, we exploit the definition of Jordan measurability. We trap the smooth set $A$ between two ``pixelated'' approximations made entirely of dyadic cubes: an inner approximation $F$ (yellow) and an outer approximation $G$ (red). Because these sets are built from orthogonal blocks, taking a 1D vertical slice at a fixed $x$ yields perfectly vertical piecewise-constant line segments. The true slice $A_x$ is rigorously trapped: $F_x \subset A_x \subset G_x$.
\end{explanation-of-steps}
\end{math-stroke}

\begin{math-stroke}[Approximation by Dyadic Cubes (Original Visual)]
\begin{center}
\begin{tikzpicture}[scale=1.5]
    % Axes
    \draw[->, thick, profgreen] (-0.5,0) -- (4,0) node[right] {$x$};
    \draw[->, thick, profblue] (0,-0.5) -- (0,3) node[above] {$y$};
    \node[above right] at (3.5, 2.5) {$A$};

    % Outer approximation (red outline)
    \draw[red, thick] (0.5,0.5) rectangle (3.5,2.5); % Bounding box

    % Inner approximation (green inside) - simplified for visual clarity
    \draw[profgreen, thick] (1.0, 0.7) -- (1.0, 1.3) -- (1.5, 1.3) -- (1.5, 0.7) -- cycle;
    \draw[profgreen, thick] (1.0, 1.7) -- (1.0, 2.3) -- (1.5, 2.3) -- (1.5, 1.7) -- cycle;
    \draw[profgreen, thick] (2.0, 0.7) -- (2.0, 1.3) -- (2.5, 1.3) -- (2.5, 0.7) -- cycle;
    \draw[profgreen, thick] (2.0, 1.7) -- (2.0, 2.3) -- (2.5, 2.3) -- (2.5, 1.7) -- cycle;

    % Outer approximation of A by union of dyadic cubes (red outline)
    \draw[red, thick, smooth] plot coordinates { (0.7,0.7) (0.7,1.5) (1.2,1.5) (1.2,2.3) (1.8,2.3) (1.8,2.7) (2.7,2.7) (2.7,1.8) (3.3,1.8) (3.3,0.7) (2.5,0.7) (2.5,0.2) (1.0,0.2) (1.0,0.7) (0.7,0.7) };

    % Inner approximation of A by union of dyadic cubes (green outline)
    \draw[profgreen, thick, smooth] plot coordinates { (0.9,0.9) (0.9,1.3) (1.3,1.3) (1.3,2.1) (1.7,2.1) (1.7,2.5) (2.3,2.5) (2.3,1.7) (2.9,1.7) (2.9,0.9) (2.0,0.9) (2.0,0.5) (1.0,0.5) (1.0,0.9) (0.9,0.9) };

    % A specific slice for x
    \draw[blue, very thick] (1.7, 0.5) -- (1.7, 2.5);
    \node[below, xshift=-0.1cm, blue] at (1.7, 0) {$x$};

    % Slices of the approximations
    \draw[densely dashed, red, thin] (1.7, 0.7) -- (1.7, 2.5);
    \draw[densely dashed, profgreen, thin] (1.7, 0.9) -- (1.7, 2.3);
\end{tikzpicture}
\end{center}
\begin{explanation-of-steps}
To prove the theorem, we rely on approximating the Jordan measurable set $A$ from both the inside and the outside using unions of small dyadic cubes. The outer approximation is shown with a red outline, and the inner approximation with a green outline. When we take a slice ($x$) of these approximations, the slices of the actual set $A$ are bounded by the slices of the approximations.
\end{explanation-of-steps}
\end{math-stroke}

\begin{spoken-clean}[00:04:08 - 00:05:00]
We will take the approximation from the inside, and then there will be the approximation from the outside. Let me describe with a picture what we are going to do in the proof, and later we will write the symbols. Since our set is Jordan measurable, we can take a very small pixel size (i.e., $2^{-p}$). We can take an inner and an outer approximation of our set by dyadic cubes, so by \qt{pixels}. The difference between these two is very small. We can make the volume of the red outer approximation minus the volume of the yellow inner approximation less than $\epsilon$ for any $\epsilon > 0$ we choose (i.e., $\mu_n(G) - \mu_n(F) < \epsilon$). Now we will see what happens when we take a vertical slice of this set. If we fix $x$ and take a slice, we see that the slice of our set $A$ contains the slice of the yellow inner approximation, and is contained within the slice of the red outer approximation (i.e., $F_x \subset A_x \subset G_x$).
\end{spoken-clean}

\begin{math-stroke}[Idea: Marginal Densities for Dyadic Approximations]
The marginal densities computed for dyadic inner and outer approximations are simply dyadic step functions.
\end{math-stroke}

\begin{spoken-clean}[00:05:00 - 00:06:00]
As you start to move $x$, you will see that whenever you move $x$ by a distance less than the size of a single pixel, you do not jump to a new cube. Therefore, the marginal densities computed for the dyadic approximations are simply dyadic step functions.
\end{spoken-clean}

\begin{math-stroke}[Definition of Dyadic Cubes in $\mathbb{R}^n$]
Given a dyadic cube (``pixel'' of size $2^{-p}$) in $\mathbb{R}^n$:
\[
Q_p^n(z) = z + 2^{-p} [0,1)^n, \quad \text{where } z = (a,b), \quad a \in \mathbb{R}^k, b \in \mathbb{R}^{n-k}
\]
The full $n$-dimensional dyadic cube can be decomposed into a product of lower-dimensional dyadic cubes:
\[
Q_p^n(z) = Q_p^k(a) \times Q_p^{n-k}(b)
\]
\end{math-stroke}

\begin{spoken-clean}[00:06:00 - 00:07:10]
Let us take the size of a pixel to be $2^{-p}$ in $\mathbb{R}^n$. We will define the dyadic cube $Q_p^n(z)$ as $z + 2^{-p} [0,1)^n$. Since there will be cubes of several dimensions, I will explicitly indicate the dimension $n$ as a sub-index. Now, if $z$ is of the form $(a,b)$ where $a \in \mathbb{R}^k$ and $b \in \mathbb{R}^{n-k}$, then in this coordinate splitting, the $n$-dimensional cube $Q_p^n(z)$ is exactly equal to the Cartesian product $Q_p^k(a) \times Q_p^{n-k}(b)$.
\end{spoken-clean}

\begin{math-stroke}[Definition of Dyadic Cubes in $\mathbb{R}^m$]
Where, in general, for $m$-dimensions:
\[
Q_p^m(\bar{z}) = \bar{z} + 2^{-p} [0,1)^m \subset \mathbb{R}^m
\]
\end{math-stroke}

\begin{spoken-clean}[00:07:10 - 00:07:37]
This is something very complicated to write, but geometrically very trivial on the picture. I am simply saying that a dyadic cube in the full space is always the product of a dyadic cube of the same size on the base axis times a dyadic cube of the same size on the vertical axis.
\end{spoken-clean}

\begin{spoken-clean}[00:07:37 - 00:08:15]
Now, given $\epsilon > 0$, let $p \in \mathbb{N}$ be such that $F$ and $G$ are dyadic sets (i.e., unions of dyadic cubes of pixel size $2^{-p}$) such that $F \subset A \subset G$, and the measure of $G$ minus the measure of $F$ (i.e., $\mu_n(G) - \mu_n(F)$) is less than $\epsilon$. This is the very definition of a Jordan measurable set. We can trap it between dyadic sets, such that the outer measure minus the inner measure is less than $\epsilon$.
\end{spoken-clean}

\begin{math-stroke}[Defining Dyadic Slices]
For all $x \in [-C,C]^k$, we define the vertical slices:
\[
G_x = \{ y \in \mathbb{R}^{n-k} \mid (x,y) \in G \} \quad \text{and} \quad F_x = \{ y \in \mathbb{R}^{n-k} \mid (x,y) \in F \}
\]
\textbf{Observation:} If $Q_p^k(a) \subset \mathbb{R}^k$ is a dyadic cube of side length $2^{-p}$, then for all $x, x' \in Q_p^k(a)$, we have:
\[
F_x = F_{x'} \quad \text{and} \quad G_x = G_{x'}
\]
\end{math-stroke}

\begin{spoken-clean}[00:08:15 - 00:09:47]
Now, for all $x \in [-C,C]^k$, we define the slices $G_x$ (i.e., the set of all $y \in \mathbb{R}^{n-k}$ such that $(x,y) \in G$), and we do similarly for $F_x$. The key observation, which I was trying to convey with the drawing, is this: if $Q_p^k(a) \subset \mathbb{R}^k$ is a dyadic base cube of pixel size $2^{-p}$, then because every cube in $\mathbb{R}^n$ is a product of cubes, $F_x = F_{x'}$ for all $x, x'$ within this base cube $Q_p^k(a)$. As long as the base point remains inside this specific cube, the vertical slices cannot change; they are perfectly constant. The same logic applies to the outer approximation $G$.
\end{spoken-clean}

\begin{math-stroke}[Defining Dyadic Step Functions]
What does it mean? It means that if we define:
\[
f(x) := \mu(F_x) \quad \text{and} \quad g(x) := \mu(G_x)
\]
Then these functions are constant in each dyadic cube $Q_p(a)$, and are zero outside of $(-C,C)^k$.
\end{math-stroke}
\end{proof}
